In this section, we comprehensively evaluate our proposed method, FedDecay, alongside other prominent federated learning techniques, using the benchmark established by \citet{pfl-bench}.
To ensure uniformity, we integrate our approach into their code base available at GitHub~\footnote{Original benchmark at  \url{https://github.com/alibaba/FederatedScope/tree/Feature/pfl_bench}}~\footnote{Our modifications at \url{https://github.com/JoeLavond/FedDecay/tree/bench}}.
We leverage this established infrastructure to experiment with identical data sets and models to facilitate a rigorous and fair comparison.

\paragraph{Data.} 
As discussed in \citet{generalization}, generalization in federated learning refers both to new users (participation gap) and new data for existing users (out-of-sample gap).
Our experiments involve holding out both users and data to study both types of generality in the following manner.
For each data set, we maintain 20 percent of users as a holdout group, simulating the challenge of generalizing to new users.
Also, we divide each user's data into distinct training, validation, and testing subsets.
Our experiments encompass the various data sets FEMNIST~\citep{femnist}, SST2~\citep{sst-92, sst-94}, and PUBMED~\citep{pubmed}.
This diversified selection of data sets empowers us to investigate performance across varied domains, including vision, text, and graph data.
We use the default data settings chosen by \citet{pfl-bench}.
\begin{itemize}
    \item FEMNIST: A 62-way handwritten character classification problem with images of size $28 \times 28$.
    This sub-sampled version of FEMNIST in \citet{pfl-bench} contains 43,400 images from 3,550 authors as clients.
    \item SST2: A sentiment classification data set containing 68,200 movie review sentences labeled with human sentiment.
    Partitioned into 50 clients using Dirichlet allocation with $\alpha = 0.4$, it enables the assessment of text-based applications.
    \item PUBMED: A graph of 19,717 nodes and 44,338 edges, this data set classifies scientific publications into three classes.
    Louvain~\citep{louvain} community partitioning partitions the graph into five users.
\end{itemize}

\paragraph{Model.} 
In our experiments, we adhere to the model architectures prescribed by \citet{pfl-bench} for consistency and comparative purposes. 
\begin{itemize}
    \item FEMNIST: We employ a Convolutional Neural Network (CNN) with a hidden size 2,048.
    The model incorporates two convolutional layers with $5 \times 5$ kernels, followed by max pooling, batch normalization, ReLU activation, and two dense layers.
    \item SST2: Our model leverages a pre-trained BERT-Tiny model~\citep{model-sst2} with 2-layer Transformer encoders and a hidden size 128.
    \item PUBMED: We employ the Graph Isomorphism Neural Network (GIN)~\citep{model-pubmed}, featuring 2-layer convolutions with batch normalization, a hidden size 64, and a dropout rate of 0.5.
\end{itemize}

\paragraph{Methods.} To establish the prowess of FedDecay, we subject it to a rigorous comparison with a range of federated learning techniques.
\begin{itemize}
    \item Federated Learning: In this category, we include FedAvg/Reptile~\citep{fedavg, reptile} and FOMAML~\citep{reptile}. 
    These methods make up the most relevant baselines to our proposal as they all can be seen, like our proposal, as various within-round learning rate scheduling.
    \item Personalized Federated Learning: These encompass Ditto~\citep{ditto}, FedBN~\citep{fedbn}, FedEM~\citep{fedem}, and pFedMe~\citep{pfedme}.
    Interpret these PFL solutions as an upper bound on what performance could be for our method as they require additional computation cost and memory.
\end{itemize}

\paragraph{Hyperparameters.} 
We extend the hyper-parameter optimization methodology outlined in the original benchmark paper.
Leveraging a grid search technique, we conduct comprehensive explorations, leveraging early termination and hyperband stopping \citep{wandb}.
Our exhaustive search covers a spectrum of hyperparameters, as detailed in Table~\ref{tab: tuning--grid}.
\begin{table*}[htbp]
    \centering
    
    \begin{adjustbox}{max width = \linewidth}
        \begin{tabular}{|l|c|c|c|}
        \hline
            Hyper-parameter & Algorithm & Data Set & Tuning Grid \\
            \hline
            Local Epochs & - & - & $\left\{ 1, 3 \right\}$ \\
            Batch Size & - & SST2 & $\left\{ 16, 32, 64 \right\}$ \\
            Learning Rate & - & - & $\left\{ 0.005, 0.01, \ldots, 1.0 \right\}$ \\
            Regularization Rate & Ditto, pFedMe & - & $\left\{ 0.05, 0.1, 0.2, 0.5, 0.9 \right\}$ \\
            Meta-learning Step & pFedMe & - & $\left\{ 1, 3 \right\}$ \\
            Mixture Number & FedEM & - & $\left\{ 3 \right\}$ \\
            Local Decay & FedDecay & - & $\left\{0, 0.2, \ldots, 1.0 \right\}$ \\
            \hline
        \end{tabular}
    \end{adjustbox}
    
    \caption{
        Hyper-parameter grid search details. If no algorithm or data set is specified, the given hyper-parameter search will be applied for all.
    }
    \label{tab: tuning--grid}
    %\vspace{-10pt}
\end{table*}


The final hyperparameters for each run are selected based on the highest average validation accuracy across users.
If runs do not terminate early, they will run for 1000, 500, and 500 epochs for FEMNIST, SST2, and PUBMED, respectively.
Other fixed configurations worth noting are 20 percent partial user participation for FEMNIST and SST2, batch size 32 for FEMNIST, and full-user participation with full-batch training for PUBMED.
To ensure a fair comparison and alignment with other methods, we maintain a fixed learning rate ($\alpha$) for FedDecay, while $\beta$ is the key hyperparameter explored.
Fine-tuning happens for one epoch with the same learning rate $\alpha$ for all methods. 
We conduct our experiments on 4 NVIDIA GeForce RTX 3090 GPUs.
See Appendix~\ref{app: decay-app-sens} to analyze FedDecay's performance under misspecified $\beta$. 